\documentclass[10pt,a4paper]{amsart}
\usepackage{amsmath,amssymb,amsthm,graphicx}
% For commutative diagrams we use Xy-pic
\usepackage[all]{xy}

% We use 2cell for 2-commutative diagrams.
\xyoption{2cell}
% Theorem environments.
%
\theoremstyle{plain}
\newtheorem{theorem}[subsection]{Theorem}
\newtheorem{proposition}[subsection]{Proposition}
\newtheorem{lemma}[subsection]{Lemma}

\theoremstyle{definition}
\newtheorem{definition}[subsection]{Definition}
\newtheorem{example}[subsection]{Example}
\newtheorem{exercise}[subsection]{Exercise}
\newtheorem{situation}[subsection]{Situation}

\theoremstyle{remark}
\newtheorem{remark}[subsection]{Remark}
\newtheorem{remarks}[subsection]{Remarks}

\numberwithin{equation}{subsection}

% Macros
%
\def\lim{\mathop{\mathrm{lim}}\nolimits}
\def\colim{\mathop{\mathrm{colim}}\nolimits}
\def\Spec{\mathop{\mathrm{Spec}}}
\def\Hom{\mathop{\mathrm{Hom}}\nolimits}
\def\Ext{\mathop{\mathrm{Ext}}\nolimits}
\def\SheafHom{\mathop{\mathcal{H}\!\mathit{om}}\nolimits}
\def\SheafExt{\mathop{\mathcal{E}\!\mathit{xt}}\nolimits}
\def\Sch{\mathit{Sch}}
\def\Mor{\mathop{\mathrm{Mor}}\nolimits}
\def\Ob{\mathop{\mathrm{Ob}}\nolimits}
\def\Sh{\mathop{\mathit{Sh}}\nolimits}
\def\NL{\mathop{N\!L}\nolimits}
\def\CH{\mathop{\mathrm{CH}}\nolimits}
\def\proetale{{pro\text{-}\acute{e}tale}}
\def\etale{{\acute{e}tale}}
\def\QCoh{\mathit{QCoh}}
\def\Ker{\mathop{\mathrm{Ker}}}
\def\Im{\mathop{\mathrm{Im}}}
\def\Coker{\mathop{\mathrm{Coker}}}
\def\Coim{\mathop{\mathrm{Coim}}}

% Boxtimes
%
\DeclareMathSymbol{\boxtimes}{\mathbin}{AMSa}{"02}

%
% Macros for moduli stacks/spaces
%
\def\QCohstack{\mathcal{QC}\!\mathit{oh}}
\def\Cohstack{\mathcal{C}\!\mathit{oh}}
\def\Spacesstack{\mathcal{S}\!\mathit{paces}}
\def\Quotfunctor{\mathrm{Quot}}
\def\Hilbfunctor{\mathrm{Hilb}}
\def\Curvesstack{\mathcal{C}\!\mathit{urves}}
\def\Polarizedstack{\mathcal{P}\!\mathit{olarized}}
\def\Complexesstack{\mathcal{C}\!\mathit{omplexes}}
% \Pic is the operator that assigns to X its picard group, usage \Pic(X)
% \Picardstack_{X/B} denotes the Picard stack of X over B
% \Picardfunctor_{X/B} denotes the Picard functor of X over B
\def\Pic{\mathop{\mathrm{Pic}}\nolimits}
\def\Picardstack{\mathcal{P}\!\mathit{ic}}
\def\Picardfunctor{\mathrm{Pic}}
\def\Deformationcategory{\mathcal{D}\!\mathit{ef}}
\title{Finite locally free sheaves with isomorphic endomorphism algebras differ by an invertible twist}
\author{The Stacks Project Authors}
\date{stacks-project:etale-cohomology.tex:local-snapshot-20261011; GNU FDL 1.2 or later}
\begin{document}
\maketitle
\noindent\textit{Editorial note.} All delivered mathematical text is the author's own, copied from the pinned Stacks Project chapter file etale-cohomology.tex; no statement or proof step was rewritten, reordered or removed.\par\medskip
\noindent\textit{Editorial note.} The proof is the author's complete argument. Two steps are described by the author as a local calculation after reducing to the affine finite free case; that reduction is stated explicitly in the counted span and is the author's own routine verification, not an omission of load-bearing mathematics.\par\medskip
\noindent\textit{Editorial note.} The label lemma-end-unique-up-to-invertible occurs in no other staged Stacks chapter, so no cross-chapter restatement had to be weighed for this card.\par\medskip
\noindent\textit{Editorial note.} The preamble is the chapter's own native preamble excerpt already used for this source: the Xy-pic setup (native preamble lines 14-18) and the native theorem declarations with the mathematical macros (native preamble lines 154-217).\par\medskip
\noindent\textit{Allowed standard prerequisites (editorial).}\begin{itemize}
\item Finite locally free sheaf of O\_S-modules of positive rank. Same chapter, the notion of a finite locally free sheaf of O\_S-modules and its rank. Supplies the sheaves F and G, and allows the local reduction to finite free modules over an affine base used in the proof.
\item Sheaf of homomorphisms and the endomorphism algebra of a finite locally free sheaf. Same chapter, SheafHom\_O(F, F) with its structure of O\_S-algebra. Supplies the algebras whose isomorphism is the hypothesis, and the ambient sheaf SheafHom(F, G) inside which the invertible sheaf L is constructed.
\item Conjugation by a local isomorphism and locally free sheaves of rank one. Same chapter, the action of a local isomorphism varphi : F -> G on endomorphism algebras by conjugation, and invertibility of a locally free sheaf of rank one. Supplies the condition selecting the generators of L, and the target property that L is invertible.
\end{itemize}\medskip
\begin{lemma}
\label{lemma-end-unique-up-to-invertible}
Let $S$ be a scheme. Let $\mathcal{F}$ and $\mathcal{G}$ be finite locally
free sheaves of $\mathcal{O}_S$-modules of positive rank. If there
exists an isomorphism
$\SheafHom_{\mathcal{O}_S}(\mathcal{F}, \mathcal{F}) \cong
\SheafHom_{\mathcal{O}_S}(\mathcal{G}, \mathcal{G})$ of
$\mathcal{O}_S$-algebras, then there exists an invertible sheaf
$\mathcal{L}$ on $S$ such that
$\mathcal{F} \otimes_{\mathcal{O}_S} \mathcal{L} \cong \mathcal{G}$
and such that this isomorphism induces the given isomorphism of
endomorphism algebras.
\end{lemma}

\begin{proof}
Fix an isomorphism
$\SheafHom_{\mathcal{O}_S}(\mathcal{F}, \mathcal{F}) \to
\SheafHom_{\mathcal{O}_S}(\mathcal{G}, \mathcal{G})$.
Consider the sheaf $\mathcal{L} \subset \SheafHom(\mathcal{F}, \mathcal{G})$
generated as an $\mathcal{O}_S$-module by the local isomorphisms
$\varphi : \mathcal{F} \to \mathcal{G}$ such that conjugation by
$\varphi$ is the given isomorphism of endomorphism algebras.
A local calculation (reducing to the case that $\mathcal{F}$ and $\mathcal{G}$
are finite free and $S$ is affine) shows that $\mathcal{L}$ is invertible.
Another local calculation shows that the evaluation map
$$
\mathcal{F} \otimes_{\mathcal{O}_S} \mathcal{L} \longrightarrow \mathcal{G}
$$
is an isomorphism.
\end{proof}

\end{document}
